% !TEX root = ../main.tex
\section{Conclusion}
\label{sec:conclusion}
We performed a Monte Carlo study on the empirical coverage of three confidence intervals based on the FKRB estimator for random coefficient discrete choice models.
We repeated the analysis for seven different distributional specifications of the random coefficients.
The analysis reveals that the empirical coverage of all three intervals on distributions that are on the boundary of the space of valid distributions is generally considerably higher than nominal.
The difference between nominal and empirical coverage is much smaller for distributions that are off the boundary, even for near-boundary distributions.
This result is in line with the theory of \citet{Andrews1999} and the observations of \citet{FoxKimRyanBajari2011}.

Among the three different intervals, the one suggested by \citet{FoxKimRyanBajari2011} has the best coverage in general, in the sense that the empirical coverage is closest to the nominal for this interval.
It is fairly consistent across different distributions, although we do see a noticeable difference between on- and off-boundary distributions.
However, this interval tends to undercover somewhat, especially for larger grids.
The sieve Wald and quasi-likelihood ratio intervals overcover, and in contrast to the FKRB interval, this is especially pronounced for larger grids.
The sieve Wald interval is the most conservative of the three; the QLR interval falls between the FKRB and sieve Wald intervals in terms of conservativeness.

All in all this thesis shows that inference on the FKRB estimator is shaky.
Researchers should be cautious when drawing any conclusions from inference in empirical applications.
With that in mind, the original FKRB interval is probably the best choice in most cases.
In situations where type I errors are costly, the QLR interval may be the better choice.

The conclusion comes with some nuances, however.
First, the results for the FKRB interval at the larger grid seem to be heavily influenced by a poorly behaved OLS estimator.
This has not been observed in other studies on the FKRB estimator, as far as the author is aware.
Whether this phenomenon is an inherent property of the unconstrained FKRB estimator or an unfortunate consequence of our DGPs would be good to know for further research.

Second, in our study we provide the correct support to the FKRB estimator.
In most applications, the support is likely misspecified.
Such misspecification can severely impact the results.
For this reason, our results may not generalize to grid-misspecified FKRB estimators.

Third, the DGPs in our simulation represent only a small fraction of all possible DGPs.
Further research might cover different distributions, different sample and grid sizes $N, R$, more (or fewer) dimensions, more (or fewer) choices $J$, or a different covariate distribution.
However, it is impossible to cover all cases in a simulation study.
The only way to address this is by developing theory.

This also brings us to the fourth nuance in our study.
None of our results are based on sound theory.
We even deliberately made some incorrect assumptions about the asymptotic distribution of test statistics on the boundary, for lack of better options.
Such theory can build on many existing studies.
\citet{Li2025Inference,Li2025The} provides theory on inference for constrained extremum estimators.
The sieve approach that we take can be a route as well, although the framework of \citet{ChenPouzo2015} is possibly excessive; \citet{ChenLintonVanKeilegom2003} might already provide much of what is needed.
The framework of \citet{ChenPouzo2015} is mostly useful for ill-posed inverse problems and functionals that are not root-$n$ estimable.
The FKRB estimator of the random coefficient model is relatively well behaved in that sense, as it is not ill-posed and is root-$n$ estimable.
However, the PSMD approach might be an interesting framework for inference on a regularized version of the FKRB estimator, such as the elastic net of \citet{HeissHetzeneckerOsterhaus2022}.

We end by stressing that our Monte Carlo study uses a limited number of repetitions ($M=250$).
One unfortunate repetition that produces an interval that does not cover the true parameter already lowers the empirical coverage by $1/250 = 0.004$.
The precision of our empirical coverage estimates is therefore limited.